\documentclass[11pt]{article}
\usepackage[margin=1in]{geometry}
\usepackage{amsmath}
\usepackage{graphicx}
\usepackage{hyperref}

\title{Heterogeneous Verification for Multi-Step Decentralized Workflows}
\author{Cory Gabrielsen (cory.eth)}
\date{\today}

\begin{document}

\maketitle

\begin{abstract}
Verifiable multi-step automation would extend blockchain guarantees beyond single-transaction execution, enabling multi-step workflows with cryptographic attestations. Smart contracts provide verification for onchain computation, but these guarantees are lost when workflows require offchain execution such as external APIs, confidential compute, or nondeterministic operations. We propose a heterogeneous verification system in which each workflow step uses the verification mechanism appropriate to its computation. The system composes cryptographic proofs for blockchain operations, hardware attestations for confidential execution, and stake-backed attestations for external interactions into unified workflow attestations. Each step produces a signed attestation that validators broadcast and stake against. Provided honest validators control sufficient stake, fraudulent attestations can be challenged via onchain fraud proofs, preserving end-to-end verifiability across heterogeneous trust assumptions. The system enables verifiable automation previously impossible under uniform trust models.
\end{abstract}

\section{Introduction}

% To be written

\section*{References}

\begin{enumerate}
\item[{[1]}] S. Nakamoto, ``Bitcoin: A Peer-to-Peer Electronic Cash System,'' 2008.
\end{enumerate}

\end{document}

